%
%============= CHAPTER BREAK =================%
%
\null%
\thispagestyle{empty}%
\clearpage%
%
\chapter{Introduction}%
%
\vspace*{5mm}
%
\justifying

\noindent A \textsc{neologism} --- a \textsc{new word} --- creates a \textit{rippling} in Time. %
It splashes forward upon its surface of porous fabric; %
echoes in waves, spoken and written, until it enters our ears and minds. %
Neologisms hide in plain sight, %
embedded among more commonplace words. %
Words that at first sounded new and creative eventually fade into the background, %
with origin points and meanings increasingly difficult to pinpoint over time. %
Massive Corpus Linguistics projects, however, are changing this. %
As more texts transform into structured data, %
the temporal topographies of human speech patterns are becoming clearer, easier to map.\endnote{%
Among recent major studies on postclassical Greek neologisms, %
see: Domazakis 2018, focusing on 2 Maccabees in particular, but with connections %
made to related neologisms in Polybius, the LXX, and apocryphal/deuterocanonical books.} 
%

For ancient Greek in particular, no project is better known %
or more useful than the Thesaurus Linguae Graecae, or the TLG, %
as it is affectionately called by its user base. %
Founded and centered at the University of California, Irvine, %
the TLG now contains well over 125 million Greek words, %
stretching from the 18th century AD all the way back to Homer and Hesiod. %
Far and away the largest pre-modern, %
publicly available digital Greek corpus in existence today, %
the TLG holds rich treasures of illumination into the history of the Greek language, %
and the evolution and transmission of human languages more generally.
Two particularly powerful features developed by the TLG facilitate %
the process of uncovering the origin points of neologisms: %
search by lemma (i.e., the dictionary or root form) and sort by date. %
Though the TLG first started in 1972, %
these features were not publicly released until late in 2006.

At present the TLG has more than 90,000 different lemmas. %
Searching and evaluating these one at a time would take years. %
Deciding which lemmas to query is often a matter of serendipity, %
dependent on the particular text or topic being researched. %
Knowing where to start and how to scope a search for neologisms might feel overwhelming.

The recent release of new, massive Greek corpora as open access data %
has decisively bent the research arc and lowered the bar to entry, 
both in terms of cost and technical feasibility. %
Open science projects such as the Diorisis (10M+ words released in 2018), %
GLAUx (now 20M+ words, first released in 2021) and Opera Graeca Adnotata %
or OGA (now 40M+ words, first released in 2024) are available to everyone. %
As more public domain texts are digitized, more agreements with copyright holders are made, %
and more new editions are released in born-digital formats, %
both GLAUx and OGA will continue to grow substantially in future releases.

These openly available Greek corpora provide root word forms, %
as well as date ranges for the composition of works. %
These ranges vary considerably in breadth and precision. %
However, with some scripting and database compiling, %
along with expert evaluation and refinement of dates, %
over 2,300 different Greek works and all of their root words %
are now available to search and to sort by date in a matter of seconds.
If one takes a database-driven approach --- ingesting, correlating, %
and querying these corpora ---, one can build a magnet to detect %
potential neologistic needles hiding in these high and heavy lexical haystacks. %
This book pinpoints these tiny treasures, %
micro-features, trace elements that together give us our clearest picture %
yet of the dates when documents in the NT were composed, edited, and/or reached stable form. %

\section{Scope, Sources, and Challenges}

\noindent{}Establishing an initial chronological line of demarcation %
proved tremendously helpful to decide %
what is more and less relevant to such an endeavor. %
For this volume, the end of Emperor Hadrianʼs reign (138 CE) supplies that line. %
So, if a Greek lemma is securely attested before or even during Hadrianʼs reign, %
then it is disqualified as a neologism in the main analysis of this book. %
This provides focus and clarity for all that follows. %
But it also squarely confronts the \textit{status quo} scholarly dates %
assumed and/or argued for New Testament texts, %
dates ranging from the mid-1st to early 2nd century CE. %
If our most reliable chronological witnesses %
point to new Greek words being coined and/or first popularized after Hadrian, %
and these words are found throughout the books of the canonical New Testament, %
then this should move the needle decisively on the outer bound dates of NT documents. %

Even with this initial line of demarcation in place, %
sifting through primary sources outside the NT to separate %
pre-138 CE lemmas from exclusively post-138 CE lemmas is still often quite challenging. %
An enormous diversity of texts by authors both known and unknown %
present myriad obstacles for scholars to navigate over, through, and around %
on their respective quests for scientific clarity. %
Many 2nd century CE Greek texts have survived in their original language, %
yet their stable and corroborated textual forms date to the 4th or 5th century. %
What may have originally been 2nd century Greek works are sometimes %
only preserved in Syriac, Coptic, or Latin translations in manuscripts dated several centuries later. %
Christian martyr Acts, for example, started to appear around the mid- to late-2nd century, %
but the versions preserved for posterity often reflect 3rd, 4th, or 5th century editions.\endnote{%
See the various introductions in Musurillo 1972, %
and more recently and especially, Moss 2013.} %
In practice, the imprisonment of Origin during the Decian persecution %
supplies a \textit{terminus ante quem} for works considered relevant to this project. %
When martyr stories with neologism candidates may have originated %
before Decius, %
they were considered in the analysis, but otherwise were ignored. %

Other works, particularly those that later became collectively %
known as the Apostolic Fathers, were often re-written and expanded, %
making precise dating of their earliest textual strata fraught.\endnote{%
Vinzent 2024; Bull 2023.} %
Citations to such texts often appear in the analysis or notes, %
but they were usually found unreliable candidates for the earliest chronological hits of a lemma.\endnote{%
For example, hits found in the corpora of letters of Ignatius of Antioch %
tend to appear in the quite late 13 letter collection, sometimes in the 7 letter collection, but almost never in Curetonʼs 3 letter collection. %
In the analysis, I label citations to the Ignatian corpora as Ign3, Ign7, and Ign13 for precisely this reason.} %
Pseudonymous works attributed to 2nd and early-3rd century authors %
--- Aristides, Euodius, Hippolytus of Rome, Irenaeus, Tertullian, %
and others,--- were deemed unreliable as providing secure samples of neologisms %
between 138 and 249 CE, and were thus generally ignored. %

Works labeled as Old Testament Apocrypha or Greek Pseudepigrapha %
are handled cautiously on a case by case basis. %
These include texts such as \textit{2--4 Maccabees}, \textit{1st--4th Enoch}, %
the \textit{Testament of Abraham}, %
\textit{Testaments of the Twelve Patriarchs}, the \textit{Ascension of Isaiah}, %
the \textit{Sibylline Oracles}, the \textit{Apocalypse of Moses}, and others. % 
Some of these texts have clear Greek witnesses and corroboration, %
but others have large lacunae in Greek manuscripts and are only well-attested in Ge'ez (Ethiopic), %
Latin, or other languages. %
Sometimes TLG results for these texts are judged as providing secure pre-138 CE attestations %
of Greek lemmas (often in conjunction with more secure authors such as Philo or Josephus), %
but at other times the results are determined to be unreliable, chronologically dubious, %
and/or the product of later Christian redaction. %
Still others, such as the 4th century CE \textit{Testament of Solomon}, %
mislabeled as 1st century in the TLG, are ignored altogether.%

Early Christian Apocrypha often appear in the main analysis or endnotes, %
but they rarely provide chronologically secure points of origin. % 
These texts call for caution in regard to possible mid- to late-2nd century provenance, %
but they can be almost entirely dismissed as providing 1st century attestations, %
despite these inaccurate labels in the TLG. %
Whether they are considered Christian Apocrypha or pre-canonical versions of Luke and the Pauline letters, %
Marcion's \textit{Evangelion} and \textit{Apostle} have their own complexities. %
Though they are noted in the following analysis, the findings do not hinge on them. %
That is to say, one could ignore Marcion's scriptures entirely and still observe the patterns documented %
regarding the emergence of new words and their earliest appropriations after 138 CE. %
The vast majority of neologisms documented here are not found in wording attested for Marcion's scriptures, %
even when attested for their canonical counterparts, %
a consistent pattern that is highly unlikely to be mere coincidence. %

Even well-known Christian authors before Irenaeus present %
challenges and confounding factors. %
Justin Martyr, for example, %
provided many of our earliest clear quotations of and citations to the canonical Gospels, %
and yet Justin may also have had a hand in editing those Gospels. %
Tatian, Justin's student and a slightly older contemporary of Irenaeus, %
also presents problems, not least because of his own heavy-handed editing %
of the four canonical Gospels into a single Gospel Harmony. %

The potential for cross-contaminations across canonical NT texts is also a major confounding factor. %
Numerous terms that could be viable neologisms in one Gospel could have spilled over into the others, %
in any number of possible directions of dependence. %
Terms that first appeared in one Pauline epistle could echo in another Pauline epistle, %
or in a Petrine epistle, whether written by the same author or a different one. %
For this reason, this book brackets the question of relationships of %
sequence and dependency among canonical NT texts. %
In essence, we treat all NT books as a single group or document class. %
Though we cite specific chapter and verse locations for words within these texts, %
when doing so we do not insist on any particular model of interdependence among them on the whole. %
To get our bearings on the overall shape of the maze of NT relationships, %
we need to set ourselves outside and above the maze, even while zooming in on micro-features %
made possible by todayʼs rigorously-annotated large data corpora. %

Loanwords present yet another methodological consideration. %
Numerous Hebrew, Aramaic, and Latin words apparently find their first known appearance in the TLG %
within canonical NT texts.\endnote{%
E.g., γέεννα / \textit{gehenna} / “valley of Hinnom” or “hell” (see also \textit{Orac.Sib.} 4.187); %
σπεκουλάτωρ / \textit{spekoulatōr} / “executioner” (Mark 6:27; cp. \textit{Act.Paul}, fr. 10, \textit{Mart.Paul.} 5); %
λεγιών / \textit{legiōn} / “legion”.%
} From a cross-cultural point of view, these could be dismissed as mere transliterations %
(i.e., sounding out a word from one language using the alphabet of another). %
But relative to the particular history of the Greek language, %
these loanwords sometimes become neologisms, taking on a life of their own in the target language. %
For a whole host of reasons, Latin loanwords generally provide clearer points of Greek origin %
than Hebrew and Aramaic loanwords do. %
Thus Latin loanwords are typically included in the following analysis, %
while Hebrew and Aramaic loanwords are not. %

Generally speaking, non-Christian Greek authors of the 2nd century CE provide far clearer historical-chronological data %
than New Testament texts. %
These authors were often well-known, public figures, %
considered in their own time as political, military, scientific, philosophical and/or religious authorities. %
To be sure, there were quite popular exceptions that were largely if not entirely pseudonymous, %
such as the \textit{Fables} of Aesop, %
the \textit{Life} or \textit{Romance} of Alexander the Great, the \textit{Corpus Hermeticum}, %
or the works of Ps-Arrian. %
These exceptions present confounding factors, i.e., uncertain dates of origin and/or versional fluidity. %
Words that appear in the \textit{Fables} of Aesop and Alexander \textit{Romance} are excluded %
as neologism candidates, because the content does not lend itself to a secure date exclusively after 138 CE. %
Other works that are obviously, fully, and clearly pseudonymous are often taken to be post-138 CE, %
but they are never used as providing secure chronological basis for a neologism. %
Among Greek novels, Chariton is taken as pre-138 CE, Achilles Tatius as dubious, %
and Xenophon of Ephesus as post-138 CE, %
yet none of these ever provide the basis for securely dating the emergence of a neologism. %
We also exclude words attested securely by a variety of late %
1st and early 2nd century CE authors such as Josephus, Dio Chrysostom, and Plutarch, %
for though their influence on NT texts is more or less well-documented,\endnote{%
On the influence of Josephus on NT texts, see esp. Mason 2003, Pervo 2006, 2009, %
and the ambitious exploration of Josephan intertextuality in Adamczewski 2013, 2014a, 2014b, 2018, 2023.} 
they all sit within the range of dates considered acceptable in NT scholarship today. %

Exclusions also apply at the morphological level, particularly for named entities (personal names and place-names). %
Prosopography (the study of personal identities, names, and family relationships) %
has additional layers of complexity to consider. %
They often appear in inscriptions (letters carved in stone), which are both numerous and diffuse. %
These data are more prone to ambiguity, not only in spelling variations %
but also frequent abbreviations in the material record. %
They are also strewn across many print volumes and various databases, %
so navigating them \textit{en masse} can be challenging. %
My initial queries turned up some 240 total named entities in the New Testament %
that are possible candidates as post-Hadrianic neologisms. %
But sifting carefully through all of these names is a massive and complex undertaking. %
Perhaps that will be the topic of a later, collaborative volume or project. %

Spelling variations for common nouns are sometimes included and sometimes not, %
depending on whether the root words have separate entries in the LSJ lexicon, %
whether the words are morphologically distinct (e.g., follow different noun or adjective declensions), %
and/or whether the variations exhibit significant temporal divergence.\endnote{%
The opening example in the chapter on Aretaeus of Cappadocia is a good case in point. %
What might appear at first glance to be a mere spelling variation %
should be considered a neologism, %
not in the sense of coining a brand new word, %
but rather as a new manifestation of a previously known and used word. %
} %

\section{Organization, Approach, and Caveats}

\noindent{}Thus, in what follows, we take a chronological, %
author-by-author approach to map neologism candidates strictly by limiting ourselves to clear, %
historically well-anchored, non-Christian uses of terms. %
Essentially, we query the literary-historical record by making no assumptions about %
when the canonical Gospels and Acts came into existence, %
treating these documents as non-factors when identifying and isolating neologisms %
in the diachronic flow of time. %
To put it differently, the starting assumption is that the Gospels and Acts, %
however creative they were literarily and theologically, %
in terms of the specific vocabulary they were using, %
were seldom inventing completely new terms and introducing them into mainstream Greek usage. %
In other words, we start from the assumption that influential Greek-writing %
cultural elites were more likely than Christian authors to pioneer neologisms, %
and that their precedents were more likely to seep into the relatively meager %
literary production of Christians in the 2nd century than for Christian authors %
to invent neologisms that were then appropriated in non-Christian elite literature. % 

Letʼs put it in more pragmatic terms. %
After building a custom database using the OGA and GLAUx data, %
we queried that database for words that 1) appear in the NT, %
2) do not appear in Philo, Josephus, and Plutarch, %
and 3) do not appear in any other Greek works prior to the 2nd century CE. %
That left us with about 750 words to examine more closely. %
About 240 of these were named entities that could be filtered quickly. %
We then ran TLG lemma searches on the remaining 500 words, sorted the results by date, %
and analyzed the first several hits outside of the NT. %
This often led down scholarly rabbit-holes about the %
date and authorship of a wide variety of texts %
that are not typically a part of NT and early Christian studies. %
But these explorations, as well as fresh close readings of neologism candidates, %
turned up many original and fascinating insights into the intricacies %
of the intertextual webs produced in the initial emergence of new lemmas. %
This approach may not have caught every possible neologism coined after 138 CE %
that appears in the New Testament, but it likely caught most of them. %

We must also acknowledge the inherent limitations of this approach. %
The Greek literary record, even with modern corpora, %
is far from complete. Tens of millions of word tokens, while numerous, %
still only capture a small fraction of the entirety %
of Greek oral and written communication between 138 and 249 CE. %
Many texts are fragmentary, many more entirely lost. %
Many neologisms may have started out as oral custom before being written down. %
Many works that may be relevant (such as the Greek loanwords found in Coptic Nag Hammadi Codices) %
are not yet available as a digital corpus to search. %
We must acknowledge the limitations, arbitrariness, %
and survival biases inherent in the data that have been preserved, digitized, and made available to query. %
Be that as it may, such limitations are endemic to all scholarship on the NT and early Christianity, %
including scholarly research that does not make use of computational queries of big data. %
Progress demands that we search as much data as are available, %
while also curating new data resources, %
as we have done by compiling a private index of Greek lemmas found in the Nag Hammadi Codices. %

So, though postclassical Greek linguistic data is limited and often fragmentary, %
it is nevertheless massive and growing. %
Modern computational resources and methods allow us to query that data %
at a scale and speed that was simply impossible for previous generations. %
When we attempt to pinpoint a neologism to a specific author, document, %
and date (or date range) of composition, %
we are not claiming that these are absolutely certain and precise conclusions. %
Some of our examples may be true neologisms, %
root words invented and first used precisely in the examples we have isolated. %
Others may only get us relatively close (temporally and/or prosopographically) %
to when and in which literary circles a lemma first emerged. %
The identification of the earliest reliable appearances of a lemma %
in effect draws a line of temporal demarcation indicating the current limit of our scientific knowledge. %
These neologistic anchors are presently as far back as we can trace the specific echoes %
cascading to us through our materiality and across time.

Caution is warranted, and yet in the pursuit of science, %
open-mindedness and probabilistic logic must remain paramount. %
We have uncovered just over forty different NT lemmas that apparently first emerged %
or were first popularized after Hadrian in non-Christian texts. %
Even if many of these are flawed and inaccurate attempts to establish origin points, %
if even half of them are accurate, then that on its own should shift the Overton window %
of date ranges scholars typically hold as reasonable for the composition, %
redaction, and/or textual stabilization of these documents. %
To us it seems far more likely that a clear majority of the neologism candidates %
we surfaced actually did precede the corresponding NT texts, %
which could establish a very high baseline of confidence for correspondingly late, %
outer bound date ranges for the composition, redaction, and/or stabilization of many NT works. %
We leave it to more capable statisticians to challenge, confirm, %
correct, and concretize our preliminary probabilistic assessments. %

The main chapters typically cover one non-Christian Greek author at a time, %
particularly any of their works likely composed after 138 CE. %
Each chapter runs through all neologism candidates for that author %
and presents them in alphabetical order by Greek spelling. %
